\section{All maximum simplices in a single halfspace clipping}\label{sec:clipping}
\begin{theorem}[Exact clipping retention]\label{thm:clipping}
Let $P$ be a full-dimensional $d$-simplex, let $H$ be a closed halfspace,
and suppose $K=P\cap H$ is full dimensional. With $n=d+1$ and the
$P$-barycentric coordinates $\lambda_i$, put
\[
 \rho_* =\max_i\left(1-\max_{x\in K}\lambda_i(x)\right).
\]
Every prescribed maximum-volume simplex $S\subset K$ satisfies
\begin{equation}\label{eq:clip-bound}
 E(K,S)\le n\rho_*.
\end{equation}
Consequently corner losses at most $\rho\ge0$ imply $E(K,S)\le n\rho$.
The coefficient $n$ is sharp for every $d\ge2$ and every $0<\rho_*<1$.
All maximum simplices and their exact excess are described below.
\end{theorem}

\subsection{The vertices and a candidate maximum}
Use affine $P$-coordinates, so $P=\conv(e_1,\ldots,e_n)$ in
$\sum_i x_i=1$. Write $H=\{f\ge0\}$ and $f_i=f(e_i)$. If no $f_i$ is
negative, $K=P$ and the result is immediate. Otherwise full dimension
implies that at least one $f_i$ is positive. Partition the indices as
\[
 I=\{i:f_i=-a_i<0\},\qquad J=\{j:f_j=b_j>0\},\qquad Z=\{z:f_z=0\}.
\]
Let $B=\max_{j\in J}b_j>0$, $J_* =\{j:b_j=B\}$, and
\[
 t_i=\frac B{a_i+B},\qquad s_i=\frac{a_i}{a_i+B}=1-t_i\quad(i\in I).
\]
The vertices of $K$ are precisely
\begin{equation}\label{eq:clip-vertices}
 e_j\ (j\in J\cup Z),\qquad
 v_{ij}=\frac{b_je_i+a_ie_j}{a_i+b_j}\quad(i\in I,\ j\in J).
\end{equation}
Indeed, a new vertex lies on an old edge crossing the cut hyperplane.
Equivalently, the active coordinate constraints leave at most two positive
coordinates, and the two-coordinate case gives this expression.

For $z\in Z$, only the vertex $e_z$ has a nonzero $z$-coordinate, so it
is mandatory in every full-dimensional vertex simplex. Its column can be
expanded out. The determinant calculation may therefore be performed on
$I\cup J$ and all $e_z$ then adjoined. Fix $j_*\in J_*$. The simplex
with every retained $e_j$ and with $v_{i j_*}$ for every $i\in I$ has
normalized determinant
\begin{equation}\label{eq:clip-D}
 D_* =\prod_{i\in I}t_i>0.
\end{equation}
Expansion in the retained standard columns proves this identity. Moreover,
\begin{equation}\label{eq:clip-loss}
 \max_K\lambda_i=t_i\ (i\in I),\qquad
 \max_K\lambda_j=1\ (j\in J\cup Z),\qquad
 \rho_* =\max_{i\in I}s_i.
\end{equation}

\subsection{Weighted rooted forests and equality}
Select $|I|+|J|$ columns from the vertices with these row indices.
Represent each $v_{ij}$ by an edge $ij$ in the bipartite graph on $I\cup J$
and each selected $e_j$ by a root mark at $j$. Include isolated row
vertices. Multiply row $i\in I$ by $a_i$ and row $j\in J$ by $-b_j$.
The edge and root columns become, respectively,
\[
 \frac{a_ib_j}{a_i+b_j}(e_i-e_j),\qquad -b_je_j.
\]
Columns in a connected component form a basis of its coordinate space
exactly when that component is a tree with one selected root. An unrooted
component on $v$ vertices has incidence rank $v-1$. A rooted component
has rank at most $v$, and full independence together with the total column
count requires exactly $v$ columns in each component. Connectivity needs
at least $v-1$ edges and full rank needs at least one root, hence there
are exactly $v-1$ edges and one root. Cycles or excess roots force
column dependence. Conversely, leaf elimination shows that an unweighted
rooted tree has determinant of absolute value one.

Orient every edge of a rooted forest toward its root. Leaf elimination
before row scaling gives the exact product
\begin{equation}\label{eq:forest-product}
 |\det W|=\prod_{\text{nonroot vertices }u}
       (\text{entry at }u\text{ of its parent-edge column}).
\end{equation}
A bad vertex $i\in I$ is never a root; if its parent is $j$, its factor
$b_j/(a_i+b_j)$ is at most $t_i$. A nonroot good vertex has factor
$a_i/(a_i+b_j)<1$. Thus every selected determinant satisfies
\[
 |\det W|\le\prod_{i\in I}t_i=D_*.
\]
Equality holds exactly when every good vertex is a root and every bad
vertex attaches directly to a member of $J_*$. All vertex maxima are
therefore
\begin{equation}\label{eq:vertex-maxima}
 \conv\bigl(\{e_j:j\in J\cup Z\}\cup\{v_{i j_i}:i\in I\}\bigr),
 \qquad j_i\in J_*,
\end{equation}
with independent choices of $j_i$. Ordered by their original bad and
retained indices, all these oriented determinants equal $D_*$. Expanding
an arbitrary simplex's determinant separately in convex combinations of
$K$-vertices proves that its absolute determinant is also at most $D_*$.

\subsection{Every prescribed maximum including nonvertex maxima}
Let $S$ be any maximum and express each of its vertices as a convex
combination of vertices of $K$. Sample the combinations independently.
The oriented determinant of $S$ is the expectation of the sampled
oriented determinant. Equality in
\[
 |\E\det W|\le\E|\det W|\le D_*
\]
forces every positive-probability tuple to be a maximum from
\eqref{eq:vertex-maxima}, with the same orientation.

If no column selected a mandatory retained vertex $e_j$ surely,
independence would give positive probability that the whole tuple omitted
it. Hence each retained vertex occupies a deterministic column. After
removing those columns, each of the disjoint groups
$G_i=\{v_{ij}:j\in J_*\}$ must appear exactly once in every tuple.
The same avoidance argument forces some column to belong surely to each
$G_i$. Distinct groups require distinct columns, accounting for all
remaining columns. Merely relabeling the vertices of this same $S$ gives
\begin{equation}\label{eq:all-maxima}
 \begin{gathered}
 S=\conv\bigl(\{e_j:j\in J\cup Z\}\cup\{w_i:i\in I\}\bigr),\\
 w_i=t_ie_i+s_iz_i,\qquad z_i\in\conv\{e_j:j\in J_*\}.
 \end{gathered}
\end{equation}
Conversely every such simplex has determinant $D_*$ and is maximum.
This includes all continuous tie cases.

\subsection{Exact excess about the same centroid}
Write $z_i=\sum_{j\in J_*}z_{ji}e_j$, with $z_{ji}\ge0$ and
$\sum_jz_{ji}=1$. The $S$-coordinates of an affine point $x$ are
\begin{equation}\label{eq:clip-inverse}
 \begin{aligned}
 \alpha_i(x)&=(1+a_i/B)x_i &&(i\in I),\\
 \alpha_j(x)&=x_j-\sum_{i\in I}(a_i/B)z_{ji}x_i &&(j\in J_*),\\
 \alpha_j(x)&=x_j &&(j\notin I\cup J_*).
 \end{aligned}
\end{equation}
For $j\in J_*$ put $B_{-j}=\max(\{b_h:h\in J,\ h\ne j\}\cup\{0\})$.
The exact largest negative coordinate magnitude is
\begin{equation}\label{eq:clip-beta}
 \beta(S)=\max_{\substack{i\in I\\j\in J_*}}
       \frac{z_{ji}a_iB_{-j}}{B(a_i+B_{-j})}.
\end{equation}
A zero $B_{-j}$ contributes zero. At $v_{ih}$ with $h\ne j$, the possibly
negative coordinate is $-z_{ji}a_ib_h/[B(a_i+b_h)]$. At $v_{ij}$ it is
$s_i(1-z_{ji})\ge0$, and retained vertices have nonnegative coordinates.
Monotonicity of $b/(a_i+b)$ gives \eqref{eq:clip-beta}; convex averaging
covers every $x\in K$. Since $B_{-j}\le B$, it follows that
$\beta(S)\le\max_i a_i/(a_i+B)=\rho_*$.

A centroid dilation by $1+u$ has all simplex coordinates at least $-u/n$.
Consequently the exact identity is
\begin{equation}\label{eq:clip-exact-E}
 E(K,S)=n\beta(S)\le n\rho_*.
\end{equation}
This proves Theorem~\ref{thm:clipping} for the prescribed simplex and its
original centroid.
If $|J_*|\ge2$, then
$\beta(S)=\max_i(s_i\max_{j\in J_*}z_{ji})$. Every vertex maximum has
$E=n\rho_*$, while balanced nonvertex maxima can have smaller excess.
If $J_* =\{j_*\}$, the maximum is unique as a set, and
\[
 \beta(S)=\max_i\frac{a_iB_2}{B(a_i+B_2)},
\]
where $B_2$ is the second-largest positive value, or zero if there is none.
If $|J|=1$, $K$ itself is a simplex and $\beta=0$.
An equal-intercept one-vertex truncation has one bad vertex and, for
$d\ge2$, at least two tied good vertices. Setting $s_i=\rho$ and taking
a vertex maximum gives $\beta=\rho$, proving the stated sharpness.
The graph in this proof concerns clipping vertices; it is not the
facet-normal spanning-two-forest construction for a $d+2$-vertex circuit.

\subsection{An arbitrary body containing an opposite facet}
\begin{theorem}[One-corner bodies]\label{thm:one-corner}
Let $P=\conv(v_0,\ldots,v_d)$ be full dimensional and let $K$ be a convex
body with $\conv(v_1,\ldots,v_d)\subset K\subset P$. Put
$t=\max_K\lambda_0>0$ and $\rho=1-t$. Every prescribed maximum simplex is
\[
 S=\conv(v_1,\ldots,v_d,q),\qquad q\in K,\quad\lambda_0(q)=t,
\]
and $E(K,S)\le(d+1)\rho$. The coefficient is sharp for $d\ge2$ and
$0<\rho<1$. No polyhedral or boundary-regularity hypothesis is imposed.
\end{theorem}
\begin{proof}
If $t=1$, the maximizing point is $v_0$, so $K=P$. Otherwise set
$L=\{x\in P:\lambda_0(x)\le t\}$. By the complete clipping classification,
its maximum volume is $t|P|$, and all maxima consist of the opposite facet
and an apex at height $t$. Compactness gives an apex $q\in K$ attaining
that height, and the entire opposite facet is in $K$. This simplex has
volume $t|P|$, so $K$ and $L$ have the same maximum volume. Every
prescribed maximum of $K$ is therefore a maximum of $L$ and has precisely
the claimed form. Its apex belongs to $K$, without any requirement that
it be an extreme point of $K$.

Write the coordinates of $q$ as $(t,q_1,\ldots,q_d)$, where
$q_j\ge0$ and $\sum_jq_j=\rho$. For $x\in K$, its $S$-coordinates are
\[
 \alpha_0(x)=\lambda_0(x)/t,\qquad
 \alpha_j(x)=\lambda_j(x)-(q_j/t)\lambda_0(x)\ge-q_j\ge-\rho.
\]
The centroid formula proves the stronger estimate
$E(K,S)\le(d+1)\max_jq_j$. Equality in the asserted bound is attained by
$L$ and an apex on one truncated edge, using a second truncated edge to
realize the negative coordinate.
\end{proof}
If the actual enclosure chain supplies a $P$ satisfying the additional
opposite-facet containment hypothesis, the theorem may be used with
$\rho=(H_de(K))^{1/m}$. That extra hypothesis is not asserted for an
arbitrary body or an unrelated enclosing simplex.
